% 04-ingenieria-requisitos.tex — Fuente: docs/informe-final/04-ingenieria-requisitos.md

\chapter{Ingeniería de requisitos}
\label{cap:ingenieria-requisitos}

Este capítulo es propio y no se fusiona con el Capítulo~\ref{cap:diseno-arquitectura}
(Diseño). Documenta el proceso completo de ingeniería de requisitos que
llevó a la especificación vigente en \texttt{v1.0.0}, siguiendo la
disciplina de ISO/IEC/IEEE~29148:2018~\cite{ISO29148} y el marco de
Pohl~\cite{POHL2010} (\S2.1).

\section{Contexto y stakeholders}

Los interesados clave del proyecto se identificaron siguiendo el modelo de
capas (\emph{stakeholder onion}), como se detalla en la Tabla~\ref{tab:stakeholders}.

Como se detalla en la Tabla~\\ref{tab:validacion-req}, 
\begin{table}[H]
\centering
\caption{Identificación de stakeholders (modelo de capas de Alexander). Fuente: \texttt{docs/requisitos/elicitacion/stakeholders.md}.}
\label{tab:stakeholders}
\begin{tabular}{@{}p{2.6cm}p{5.2cm}p{6.2cm}@{}}
\toprule
\textbf{Capa} & \textbf{Interesados} & \textbf{Interés principal} \\
\midrule
Núcleo (equipo operador) & Los cuatro integrantes del equipo de desarrollo &
  Entregar un sistema funcional, evaluable y mantenible dentro del plazo
  académico \\
Usuarios directos & Creador de contenido, Cliente registrado, Administrador
  (roles definidos en \texttt{docs/requisitos/SRS.md}~\S2.3) &
  Contratar/ofrecer servicios con garantías de pago, comunicación y
  verificación de identidad \\
Usuarios indirectos & Visitante anónimo (explora el catálogo sin
  registrarse) & Evaluar la plataforma antes de registrarse \\
Patrocinador / regulador académico & Docente-director del PFC & Verificar
  que el sistema cumple los criterios de la rúbrica de cada entrega \\
Proveedores externos & PayPal (pagos), servicio de IA de verificación,
  proveedor de almacenamiento de objetos (Azure Blob / local) &
  Disponibilidad y estabilidad de la integración \\
Competencia / referencia de mercado & Plataformas genéricas de freelance y
  redes sociales usadas hoy de forma fragmentada
  (Capítulo~\ref{cap:introduccion}) & No son stakeholders directos, pero
  delimitan el estándar de usabilidad y seguridad esperado \\
\bottomrule
\end{tabular}
\end{table}

\section{Técnicas de elicitación aplicadas}

Las técnicas de elicitación aplicadas se documentan en
\texttt{docs/requisitos/elicitacion/README.md}, que declara tanto las que
se usaron como las que no, junto con la evidencia conservada de cada una.

Se aplicaron dos técnicas con evidencia archivada. La primera, el
\textbf{análisis de documentos}: el corpus original de requisitos
(\texttt{docs/requisitos/historico/entrega-1a.pdf}, semana del 4 de junio
de 2026, referenciado en \texttt{docs/requisitos/SRS.md}~\S1.4) es la
fuente de la que derivan los
37 requisitos vigentes. La segunda, el \textbf{prototipado de baja
fidelidad}: seis wireframes archivados en
\texttt{docs/diagramas/wireframes/} fijaron el alcance de las pantallas
principales antes de implementarlas. Se aplicó además la
\textbf{discusión interna del equipo} sobre el dominio de la economía
creativa, de la que no se conservaron actas ni notas de sesión.

\textbf{Limitación metodológica declarada:} no hubo entrevistas
semi-estructuradas, observación de usuarios, cuestionarios ni workshops
con stakeholders externos. El corpus se elicitó \emph{sin usuarios
finales reales}: refleja la comprensión del dominio que tiene el propio
equipo y la comparación con plataformas del sector, no necesidades
declaradas por Creadores o Clientes. Deliberadamente no se reconstruyeron
notas ni guiones a posteriori para rellenar el expediente: fabricar esa
evidencia después del hecho invalidaría justamente lo que la trazabilidad
pretende acreditar.

Esta limitación tiene dos consecuencias que se recogen en el
Capítulo~\ref{cap:validez}. Por un lado, un \emph{sesgo de constructo}:
el estudio SUS con 16 participantes mide la usabilidad del sistema
construido, no la pertinencia de los requisitos que lo originaron. Por
otro, la tasa de estabilidad de requisitos del 100\,\% (\S7.4 del SRS) es
coherente con un corpus que nunca se renegoció con un cliente externo, y
por eso no debe leerse como un mérito de la especificación.

\section{Requisitos funcionales y no funcionales: categorización y priorización}

El corpus completo está en \texttt{docs/requisitos/SRS.md} (37
requisitos: 23 funcionales \texttt{REQ-F-001}--\texttt{REQ-F-023}, 14 no
funcionales \texttt{REQ-NF-001}--\texttt{REQ-NF-014}), priorizado con
\gls{moscow}:

\begin{table}[H]
\centering
\caption{Distribución de requisitos por prioridad MoSCoW}
\label{tab:moscow}
\begin{tabular}{@{}lrr@{}}
\toprule
\textbf{Prioridad} & \textbf{Cantidad} & \textbf{\%} \\
\midrule
Must  & 29 & 78.4\,\% \\
Should &  7 & 18.9\,\% \\
Could &  1 &  2.7\,\% \\
Won't &  0 &  0\,\% \\
\bottomrule
\end{tabular}
\end{table}

\textbf{Nota de conformidad C9 (INCOSE~v4), declarada sin maquillaje:} los
37 requisitos están redactados en prosa descriptiva estructurada
(\emph{rationale}, prioridad, criterio de aceptación, verificación,
estado), \textbf{no} en el patrón sintáctico formal \emph{[condición]
[sujeto] shall [acción] [objeto] [restricción]} que exige explícitamente
ISO/IEC/IEEE~29148:2018 y las 42 reglas de
INCOSE~v4~\cite{ISO29148,INCOSE2023}.
\texttt{docs/checklists/incose-checklist.md} documenta este hallazgo (C9,
no cumplida) como el de mayor esfuerzo pendiente entre los 15 criterios
evaluados --- reescribir el corpus completo al patrón formal, no un
ajuste puntual.

\section{Modelado de casos de uso y de historias de usuario}

\begin{itemize}
  \item \textbf{Historias de usuario} (formato Connextra, criterios
    INVEST, \emph{acceptance criteria} en Gherkin):
    \texttt{docs/requisitos/historias/} (HU-01 a HU-23), cada una trazada
    a al menos un \texttt{REQ-F-NNN} y a una prueba automatizada de
    aceptación.
  \item \textbf{Casos de uso} (plantilla de Cockburn~\cite{COCKBURN2000},
    niveles 1--4): \texttt{docs/requisitos/casos-de-uso/} (CU-01 a
    CU-23), cada uno trazado a un flujo y a al menos una prueba de
    integración.
  \item \textbf{Diagrama de casos de uso UML de alto nivel:}
    Se realizó un análisis estructurado de los casos de uso para representar
    el comportamiento funcional del sistema frente a sus distintos perfiles.
    A nivel general, se definió un modelo que ilustra la interacción de los
    principales actores (Administrador, Creador, Cliente, Moderador, Soporte
    y Auditor) con los siete grandes subsistemas o macro-procesos de Artisync.
    Adicionalmente, el análisis se desglosó en catorce diagramas de casos de
    uso de alto nivel (dos por cada módulo), los cuales detallan operaciones
    críticas como la autenticación 2FA, la validación de portafolios mediante
    Inteligencia Artificial, la retención de fondos en el sistema Escrow y
    la moderación comunitaria. Todos los modelos, construidos bajo sintaxis
    UML/Mermaid, se encuentran anexados y documentados en el archivo
    \texttt{docs/informe-final/diagramas\_casos\_de\_uso.md}.
\end{itemize}

\subsection{Análisis INVEST de Historias de Usuario}

Las 23 historias de usuario formuladas para Artisync han sido evaluadas rigurosamente y cumplen con los seis atributos del modelo INVEST (Independent, Negotiable, Valuable, Estimable, Small, Testable). El análisis detallado de cada una se encuentra archivado en \texttt{docs/requisitos/historias/}. A modo de ilustración, se reproduce la evaluación de la historia \textbf{HU-01} (Iniciar sesión):

\begin{quote}
\textbf{INVEST:} Independiente del resto del flujo de autenticación; negociable en los campos exactos del formulario; valiosa (es la puerta de entrada al sistema); estimable y pequeña (un solo formulario + persistencia); testable mediante el criterio de aceptación.
\end{quote}

Este mismo rigor analítico fue aplicado individualmente a la totalidad del corpus de historias (HU-01 a HU-23), garantizando que todas representen unidades de trabajo claras, priorizables y comprobables.

\section{Validación de requisitos}

El procedimiento de validación aplicado combina dos mecanismos
verificables en el repositorio:

\begin{enumerate}
  \item \textbf{Prueba automatizada de aceptación:} cada requisito
    \emph{Must}/\emph{Should} tiene un campo \emph{Verificación} explícito
    en el SRS (Test/análisis/demostración) y una prueba automatizada
    referenciada en \texttt{docs/trazabilidad/matriz.csv} (columna
    \texttt{prueba\_automatizada}).
  \item \textbf{Estado verificado en la matriz de trazabilidad:} un
    requisito se marca \texttt{verificado} solo cuando su prueba
    automatizada está en verde contra el código en ejecución, no por
    declaración del equipo --- la re-auditoría de
    \texttt{docs/observaciones/INFORME-BRECHAS-ENTREGA-FINAL.md} confirmó
    esta disciplina detectando y corrigiendo el único caso donde el
    estado no coincidía con la evidencia real
    (\texttt{REQ-F-006}/\texttt{007}, ver
    \texttt{docs/checklists/incose-checklist.md}, hallazgo C8).
\end{enumerate}

\textbf{Resultado de la validación final (2026-08-17, contra \texttt{matriz.csv}):}

\begin{table}[H]
\centering
\caption{Estado de validación de los 37 requisitos}
\label{tab:validacion-req}
\begin{tabular}{@{}lrr@{}}
\toprule
\textbf{Estado} & \textbf{Requisitos} & \textbf{\%} \\
\midrule
Verificado                          & 25 & 67.6\,\% \\
Implementado (no verificado aún)    &  8 & 21.6\,\% \\
Pendiente                           &  4 & 10.8\,\% \\
\bottomrule
\end{tabular}
\vspace{0.3em}

{\footnotesize Los 4 requisitos \emph{Pendiente}: \texttt{REQ-F-009},
\texttt{REQ-F-010}, \texttt{REQ-NF-005}, \texttt{REQ-NF-006}.}
\end{table}

Los 29 requisitos \emph{Must} están en estado implementado o verificado
(ninguno pendiente); los 4 \emph{Pendiente} corresponden a requisitos
\emph{Should}/\emph{Could} --- funciones sociales (seguidores,
comentarios) y una medición de latencia WebSocket no ejecutada
formalmente, ambos justificables como trabajo futuro
(Capítulo~\ref{cap:trabajo-futuro}) sin bloquear el cierre de los
requisitos obligatorios.

\section{Gestión del cambio de requisitos}

\texttt{docs/requisitos/CHANGELOG-REQ.md} registra 2 versiones
(\texttt{v0.3.0} $\to$ \texttt{v0.9.0-rc}: renombrado de identificadores
\texttt{RF-NN}$\to$\texttt{REQ-F-0NN}, sin cambio semántico). \textbf{Tasa
de estabilidad declarada con salvedad} (ver
\texttt{docs/checklists/incose-checklist.md}, métricas de calidad):
nominalmente 100\,\% (0 modificaciones semánticas registradas), pero
\texttt{matriz.csv} refleja cambios de estado reales (ej.
\texttt{REQ-F-023} pasó de \texttt{pendiente} a \texttt{verificado}) sin
la entrada correspondiente en el changelog --- la tasa de 100\,\%
probablemente sobreestima la estabilidad real por subregistro, no por
ausencia genuina de cambios. Corregir este subregistro es una acción de
bajo esfuerzo pendiente para el cierre.

\section{Métricas de calidad de requisitos}

Ver el detalle completo, con evidencia citada característica por
característica, en \texttt{docs/checklists/incose-checklist.md}. Resumen:

\begin{itemize}
  \item \textbf{Número total de requisitos:} 37 (23 funcionales, 14 no
    funcionales).
  \item \textbf{Distribución por tipo:} funcionales 62.2\,\% (23/37), no
    funcionales 37.8\,\% (14/37).
  \item \textbf{Distribución por prioridad MoSCoW:} ver
    Tabla~\ref{tab:moscow}.
  \item \textbf{Porcentaje verificado:} 67.6\,\% (25/37); implementado o
    verificado (no pendiente): 89.2\,\% (33/37).
  \item \textbf{Cobertura de trazabilidad hacia código y pruebas:}
    100\,\% de los 37 requisitos tiene fila en \texttt{matriz.csv} con
    módulo de código y prueba automatizada referenciados (columnas
    \texttt{modulo\_codigo}, \texttt{prueba\_automatizada}), confirmado
    por \texttt{scripts/validate-traceability.sh} en CI.
  \item \textbf{Características INCOSE v4 cumplidas:} 9 de 15 (C1, C2,
    C4, C6, C7, C11, C12, C13, C14); 5 parciales con evidencia concreta
    (C3, C8, C10, C15, C2 parcial); 2 no cumplidas (C5, C9) --- ver
    detalle y plan de acción en el checklist referenciado.
\end{itemize}
